\documentclass[main.tex]{subfiles}


\begin{document}

%from pagenumber to up
% \phantomsection 
% \hypertarget{toc}{}

 

 

 
   

\section{Электрический ток}


\textbf{Электрический ток}~--- это направленное движение заряженных частиц. На рис.\,\ref{fig:p117} изображен ток в теле.

    \begin{figure}[H]
    \centering
\includegraphics[]{el/21I_pic.pdf}%
\caption{Электрический ток}
\label{fig:p117}
\end{figure}

Положительные заряды (красные шары) упорядоченно двигаются из одной области тела (светло"=оранжевая фигура) в другую. Через \emph{поперечное сечение} (серая плоскость) тела при протекании тока переносится заряд (равный сумме зарядов заряженных частиц, прошедших поперечное сечение за данное время). \emph{Происходит перенос заряда из одной области пространства в другую.}

% Электрический ток называют \emph{постоянным}, если за равные промежутки 

% Электрический ток может возникать в твердых телах, жидкостях и газах.

\textbf{Сила тока} ($I$ [А])~--- это характеристика тока, показывающая быстроту\footnote{Более точно, сила тока есть скорость переноса заряда $(I=\Delta q/\Delta t)$ или производная заряда по времени $(I=q')$. Важно заметить, что сила тока \emph{не} есть скорость упорядоченного движения самих частиц, образующих ток.} переноса заряда:
\begin{equation}
    I=\frac{q}{t},
\end{equation}
где $q$~--- заряд, перенесенный через поперечное сечение за время~$t$.

Ток называют \emph{постоянным}, если его сила тока остается постоянной.

Обычно скорость направленного движения заряженных частиц, образующих ток, сравнительно невелика (доли миллиметра в секунду). 

Вот \textbf{связь силы тока со скоростью} направленного движения заряженных частиц:
\begin{equation}
    I=nevS,
\end{equation}
где $n$~--- концентрация частиц, $e$~--- заряд одной частицы, $v$~--- скорость направленного движения частиц, $S$~--- площадь поперечного сечения.

Для возникновения тока необходимо выполнение следующих условий.
\begin{enumerate}
    \item \emph{Наличие свободных заряженных частиц.} Среды, содержащие такие частицы, называют \emph{проводниками}.

    \item \emph{Наличие электрического поля внутри проводника.} Иначе говоря, между концами проводника нужно обеспечить напряжение.

    Если между концами проводника создают \emph{постоянное напряжение}, то в проводнике устанавливается \emph{постоянный ток}. Постоянное напряжение поддерживают с помощью \emph{источника тока} (например, батарейки).
\end{enumerate}

Ток в \emph{металлах} обусловлен движением свободных \emph{электронов}.









 
\end{document}